\documentclass[10pt,a4paper]{amsart}
\usepackage[margin=19mm]{geometry}
\usepackage{amsmath,amssymb,mathtools}
\usepackage[scr=rsfs,cal=euler]{mathalfa}
\usepackage{xfrac}
\usepackage[unicode,hidelinks,bookmarks=false]{hyperref}
\newcommand{\ie}{i.e.\ }
\newcommand{\cf}{cf.\ }
\newcommand{\resp}{respectively\ }
\newcommand{\Subsection}{\subsection}
\providecommand{\nobreakdash}{\nobreak\hbox{-}}
\setlength{\emergencystretch}{2em}
\multlinegap0pt
\usepackage{tikz}
\usetikzlibrary{cd}
\usepackage{tikz}
\usepackage{xcolor}
\usepackage{amssymb}
\usepackage[shortlabels]{enumitem}
\usepackage{bm}
\newtheorem{lemma}{Lemma}
\newtheorem{proposition}{Proposition}
\newcommand{\coker}{\operatorname{coker}}
\title{Motivic Homotopy Groups of Spheres and Free Summands of Stably Free Modules}
\author{Sebastian Gant, Ben Williams}
\date{Source: arXiv:2510.10033v3 (CC BY 4.0)}
\begin{document}
\maketitle

\noindent\emph{Source and licence.} Motivic Homotopy Groups of Spheres and Free Summands of Stably Free Modules. Version arXiv:2510.10033v3 (CC BY 4.0), distributed by arXiv under the Creative Commons Attribution 4.0 International licence, http://creativecommons.org/licenses/by/4.0/, as recorded in the arXiv licence metadata for this exact version and shown on the abstract page https://arxiv.org/abs/2510.10033v3. Full licence notice: \texttt{licenses/arxiv-2510.10033-CC-BY-4.0.txt}.\par\smallskip

\noindent\textit{Editorial scope.} Counted unit: the article's proposition pr:5lemmaDiv (source0.tex lines 1426--1436). The article's complete original proof of it is delivered with it (source0.tex lines 1437--1466, one \texttt{proof} environment of 30 lines). No mathematics is written, repaired or re-derived: the statement and the proof are the article's own lines, in the article's own notation, and no retained line is substituted or re-flowed.  Retained uncounted as the source-local context the proof needs: lines 1417--1421.  Nothing was omitted from any retained span, so the package carries no omission record.  No retained line contains a citation, so the wrapper carries no bibliography and no bibliography entry.  The retained lines contain the article's own diagram code, which the wrapper typesets with the standard tikz packages; no image file is delivered and the package declares no asset dependency.  Editorial formatting only, in addition: an AMS \texttt{amsart} preamble that restates the article's own declarations and macros used by the retained lines, namely the environments \texttt{lemma}, \texttt{proposition} on separate counters, together with the standard packages those lines need.  The retained environments print in this document as Lemma 1, Proposition 1 in order of appearance, whereas the article prints the same items with the numbering of its own full document.  The complete native source of the article as distributed in the arXiv e-print for this exact version is bundled unchanged as \texttt{sources/186849/source0.tex}.
\begin{lemma} \label{lem:2of3divisibility}
  If
  \[ 0 \to A \to B \to C \to 0 \]
  is a short exact sequence of abelian groups with the property that two of the groups $A,B,C$ are uniquely divisible, then so is the third.
\end{lemma}

\begin{proposition}\label{pr:5lemmaDiv}
  Consider a commutative diagram of abelian groups
  \begin{equation}\label{eq:5lemma}
    \begin{tikzcd}
      A_1 \arrow[d,two heads,"\phi_1"] \rar & A_2 \arrow[d,two heads,"\phi_2"] \rar & A_3 \arrow[d,"\phi_3"] \rar & A_4 \arrow[d,two heads,"\phi_4"] \rar & A_5 \arrow[d,two heads,"\phi_5"] \\
      B_1 \rar & B_2 \rar & B_3 \rar & B_4 \rar & B_5,
    \end{tikzcd}
  \end{equation}
  in which the rows are exact sequences.
  Suppose that $\phi_1, \phi_2, \phi_3, \phi_4$ are surjective and their kernels are uniquely divisible. Suppose further that $B_3$ is torsion. Then $\phi_3$ is surjective and $\ker(\phi_3)$ is uniquely divisible.
\end{proposition}

\begin{proof}
  Treat diagram \eqref{eq:5lemma} as a double complex where $A_i$ is in degree $(i,0)$ and $B_i$ is in degree $(i,1)$. There are two spectral sequences calculating the homology of the total complex: for each of these, the nonzero groups on the $E_0$-page coincide with the groups in \eqref{eq:5lemma}. In one, the $d_0$ differentials are the horizontal arrows in \eqref{eq:5lemma}, and since the rows are exact, we see that the homology of the total complex vanishes in total degrees $3$ and $4$.

  The other spectral sequence has as $E_0$-page the following diagram
   \begin{equation*}
    \begin{tikzcd}
      A_1 \dar{\phi_1}  & A_2 \dar{\phi_2}  & A_3 \dar{\phi_3}  & A_4 \dar{\phi_4}  & A_5 \dar{\phi_5} \\
      B_1  & B_2  & B_3  & B_4 &  B_5,
    \end{tikzcd}
  \end{equation*}
  so that the $E_1$-page is
  \begin{equation}
    \begin{tikzcd}\label{eq:kernels}
      \ker(\phi_1) \rar & \ker(\phi_2) \rar & \ker(\phi_3) \rar & \ker(\phi_4) \rar & \ker(\phi_5) \\
      0 \rar & 0 \rar & \coker(\phi_3) \rar & 0 \rar  & 0.
    \end{tikzcd}
  \end{equation}
  On the $E_2$-page and later, there are no nonzero differentials arriving at or emanating from the $(3,0)$-group, which must be $0$ by the $E_\infty$-page, so that we deduce that the upper sequence in \eqref{eq:kernels} is exact at $\ker(\phi_3)$. The terms in this sequence, other than $\ker(\phi_3)$, are uniquely divisible by hypothesis, so that the five-lemma implies that $\ker(\phi_3)$ is also uniquely divisible.

  The $E_2$-page takes the form
  \[
  \begin{tikzcd}\label{eq:E2}
     \ast  & \ast   & 0  & 0  & Q \\
      0 \arrow[urr] & 0 \arrow[urr] & \coker(\phi_3) \arrow[urr, hook]  & 0   & 0.
    \end{tikzcd} 
  \]
  where $Q$ is the cokernel of $\ker(\phi_4) \to \ker(\phi_5)$ and $\ast$ denotes some unspecified abelian groups. For dimensional reasons, there are no further nonzero differentials, from which it follows that $\coker(\phi_3) \to Q$ must be an injection: the abutment of the sequence is $0$ in total degree $4$.

  Since $\ker(\phi_4)$ is uniquely divisible, its image in $\ker(\phi_5)$ is divisible, and therefore uniquely divisible. It follows from Lemma \ref{lem:2of3divisibility} that $Q$ is uniquely divisible. In particular, the map $\coker(\phi_3) \to Q$, which has a torsion domain, must be $0$. Since it is also an injection, we deduce that $\coker(\phi_3)$ is $0$ as required.
\end{proof}

\end{document}
